\subsection{A residue obstruction and a stronger conductor gap}

The complete argument and exact primary-source formulas passed an
independent root audit.

\begin{lemma}[F13g: a negative Swan intersection forces a $p$-divisible coefficient]
Let $X\to\operatorname{Spec}A$ be an SNC resolution in the
equal-characteristic complete-local surface setting of F10, and
let $S=\sum_i s_iC_i$ be the Swan divisor of a finite
$p$-primary character on the punctured spectrum. If $s_i>0$
and $p\nmid s_i$, then $S\cdot C_i\geq0$. Consequently, if
$S\ne0$, there is a component $C_i$ such that
\[
                 S\cdot C_i<0\quad\hbox{and}\quad p\mid s_i.
 \tag{F13.11}
\]
In particular $S\ne0$ implies $\max_i s_i\geq p$.
\end{lemma}
\begin{proof}
We first check the local residue explicitly. At the completed
DVR of $C_i$, put $n=s_i$ and choose a uniformizer $z$.
The character has a Witt representative
$a=(a_{m-1},\ldots,a_0)$ in the filtration
$p^jv(a_j)\geq-n$. For $j\geq1$, the assumption $p\nmid n$
implies
\[
 p^jv(a_j)\geq-p^j\lfloor n/p^j\rfloor\geq-(n-1).
\]
The refined Swan conductor is represented by
\[
                  -\sum_{j=0}^{m-1}a_j^{p^j-1}\,da_j
 \tag{F13.12}
\]
modulo logarithmic pole order at most $n-1$; these coordinate
conventions and the formula are displayed in
\cite[Definition~2.1 and (2.5), pp.~8--9]{Yatagawa2017Cycle}.
All terms with $j\geq1$ have pole order at most $n-1$.
The remaining coordinate must have pole exactly $n$, since
otherwise the character would have Swan conductor at most
$n-1$. Write $a_0=uz^{-n}$ with $u$ a unit. The logarithmic
residue of $-da_0$ in pole degree $n$ is
$n\bar u$, which is nonzero.

The global refined-Swan section of F10 maps, under the residue
sequence (F10.4) tensored by $\mathcal O_X(S)|_{C_i}$, to a
section of $\mathcal O_X(S)|_{C_i}$. Its generic value is the
nonzero residue just calculated; this generic compatibility is
Krishna--Majumder, Corollary~9.6, and the global map is
Corollary~9.7. Their Lemma~9.11 identifies this local refined
conductor with Kato's refined Swan conductor used in (F13.12)
\cite[Corollaries~9.6--9.7, pp.~55--56, and Lemma~9.11,
p.~57]{KrishnaMajumder2026}.
Thus this line bundle on the proper smooth curve $C_i$ has a
nonzero section. Its degree $S\cdot C_i$ is nonnegative.

Suppose $S\ne0$. Negative definiteness gives
\[
                  0>S^2=\sum_i s_i(S\cdot C_i).
\]
There is therefore an $i$ with $s_i>0$ and $S\cdot C_i<0$.
The first assertion forces $p\mid s_i$, proving (F13.11) and
the final conclusion.
\end{proof}

\begin{corollary}[F13h: improved unramified exponent]
In F13b and F13c, one may choose any integer $e\geq0$ with
\[
                           p^{e+1}>T_0=\max_i t_i.
 \tag{F13.13}
\]
For every rank-one $p$-primary character on every finite \'etale
resolution cover, its $p^e$ multiple then extends unramified
over that cover's resolution. Hence F13c's bounded-order twist
decomposition holds with this improved $e$.
Likewise the $p$-primary inertia kernel in F10 is killed by
$p^e$ whenever $p^{e+1}>\max_i t_i$ for the resolution used
there.
\end{corollary}
\begin{proof}
The Witt restriction step (F10.6) gives every coefficient of the
Swan divisor of $p^e\chi$ at most
$\lfloor T_0/p^e\rfloor<p$. Apply F13g on the resolution of
the complete local normalization underlying the cover. Its
exceptional matrix is negative definite, so that Swan divisor
must be zero. Purity extends the character to the resolution.
The divisible unramified character group from F13c gives the
same twist decomposition with exponent $p^e$.
For F10's abelian inertia kernel, all characters of its
$p$-primary part are killed after multiplying by $p^e$;
character separation gives the last assertion exactly as in
F10.
\end{proof}

\paragraph{Numerical consequence.}
F13f's sufficient hypothesis can therefore be weakened from
$t_i<1$ to $t_i<p$ for every $i$. Under the latter hypothesis,
F13g forces every rank-one inducing character's $p$-primary
part to be unramified; the rest of the proof of F13f is
unchanged. Neither inequality is asserted to hold for an
arbitrary normal surface after the fixed cover.
